% Kapitel 6 - Simulator-Based QA and Dataset Generation
\section{Simulator-Based QA and Dataset Generation}
\label{sec:simqa}

This chapter describes the simulator-dependent QA layer and the dataset generation procedure used for perception experiments. Simulator execution is treated as an empirical validation layer; failures are logged and preserved rather than silently repaired \cite{carla_opendrive_standalone}.

\subsection{Pipeline Validation Framework}
The pipeline applies an automated gate set before simulator execution so malformed OpenDRIVE artifacts do not propagate into experiments. These checks cover schema validity, geometric continuity, lane-graph consistency, elevation plausibility, and basic drivability. For the final governed thesis path, the domain-gap gates relevant to the reported results passed, while some finer continuity checks remained sensitive to CARLA's known \texttt{paramPoly3} approximation behavior rather than to thesis-critical structural invalidity.

The DEM coverage gate is also interpreted conservatively. Coverage reached 92.7\% after the CRS reprojection fix, which shows that the raster spatially overlaps the generated Ingolstadt road network. This does not by itself prove that non-flat elevation survived into the governed final analysis artifact: the final \texttt{08\_final\_structural\_gap.xodr} remained flat, so the reported thesis metrics are planar-only.

During hardening, the simulator QA layer was made more failure-aware through UTF-8-safe logging, clearer DEM diagnostics, deterministic spawn retries, and explicit adjacency checks for tiled validation. These controls matter methodologically because they convert silent runtime failures into auditable outcomes instead of allowing them to contaminate the reported evidence.

\subsection{Determinism evaluation protocol}
Pipeline determinism is evaluated by running the full pipeline twice with the same settings snapshot and then comparing artifacts with an external audit script. The audit checks:
\begin{itemize}
\item the settings snapshot hash,
\item the presence/absence of required artifacts,
\item hash equality of core products (final tiles, metadata, reports),
\item structured diffs of intermediate OpenDRIVE files where relevant.
\end{itemize}
This design makes the audit causally independent from the pipeline.

\paragraph{Geometric determinism versus byte-level determinism.}
Repeated conversion runs on the same fixed OSM cutout produced identical planView geometry hashes, while the raw OpenDRIVE XML files differed at the byte level (e.g., MD5). This indicates that the pipeline is \emph{geometrically deterministic} but not necessarily \emph{byte-stable}. The most common causes are non-semantic serialization variation, such as different ordering of XML elements (roads, lane sections, objects), regenerated identifiers, run-specific metadata (timestamps, run directories), or floating-point formatting differences (e.g., \texttt{0.5} vs. \texttt{0.500000}). Consequently, determinism is reported in three levels: (i) \emph{geometric determinism} (planView hash equality), (ii) \emph{structural determinism} (tile metadata and adjacency stability), and (iii) \emph{serialization determinism} (raw file hash equality). Optional canonicalization (stable sorting and formatting) can be used to enforce byte stability, but is not required for domain-gap metrics as long as the driving-relevant geometry is identical.




\subsection{CARLA environment setup (v0.9.16)}
CARLA is used in a client--server setup in which the simulator provides physics and rendering while the pipeline connects through the Python API \cite{dosovitskiy2017carla}. The packaged 0.9.16 release is used throughout because it is the stable version against which the governed evidence and runtime notes were produced \cite{carla_docs_downloads_0916}. Readiness is validated in two steps: TCP availability alone is not treated as sufficient, and actor setup begins only after an RPC-level world query succeeds \cite{carla_docs_quickstart}.

For fair comparison across maps, the simulator is run in synchronous mode with a fixed time step and explicit \texttt{world.tick()} control \cite{carla_docs_synchronous}. Generated maps are loaded primarily through CARLA's OpenDRIVE standalone workflow, which allows an XODR file to be imported directly without packaging a full Unreal map \cite{carla_opendrive_standalone}. This keeps the simulator-dependent layer lightweight and reproducible, even though richer static scene realization would require the separate cooked-map asset path.

Where large maps would otherwise exceed local GPU limits, \texttt{no\_rendering\_mode} can be used to stabilize structural QA and non-visual data collection. Any perceptual use of headless execution is reported conservatively because rendered image appearance may depend on that choice \cite{carla_python_api}.
\subsection{CARLA validation: lightweight vs.\ full-map checks}
Full-city ingestion remains unstable at scale, so simulator validation is performed primarily at tile level: the XODR is loaded via standalone mode, waypoints are generated as a navigability proxy, ego spawn is attempted, and a short bounded smoke run records the result \cite{carla_opendrive_standalone}. Tiles are then classified by failure signature (for example, load failure, waypoint failure, spawn failure, or stuck state), which yields a failure taxonomy rather than a binary pass/fail narrative.

The runtime layer was also hardened so that spawn outcomes are reported truthfully and reproducibly. Worker processes emit structured tile results, deterministic spawn retries reduce random false negatives, and adjacency information is normalized before seam analysis or streaming decisions are made. For qualitative capture, the development stress-test tile was fixed in advance by a documented loadability and road-density heuristic; this choice supports simulator QA only and is not used as primary structural evidence.


\subsection{Qualitative visual inspection of captured scenes}
To complement automated gates, representative captured scenes are inspected visually for sensor-orientation plausibility and visible scene context. This step is not used as a quantitative metric; it serves as a qualitative sanity check alongside the rig-validation reports.

\begin{figure}[ht]
\centering
\includegraphics[width=0.85\linewidth]{../images/fig_camera_rig_proof.png}
\caption{Representative RGB frames from the \texttt{camera\_rig\_fix\_v4} Town10HD\_Opt validation run. The views support the named-axis cosine verification by showing forward-facing exterior scenes for the front cameras and rearward exterior scenes for the back cameras; this figure is qualitative and not used as a quantitative metric.}
\label{fig:sim_objects}
\end{figure}

\begin{figure}[ht]
\centering
\includegraphics[width=0.54\linewidth]{../images/fig_lidar_runtime_bev.png}
\caption{Top-down LiDAR/runtime sanity visualization from the same controlled Town10HD\_Opt rig-validation family. It was used to confirm internally consistent actor placement and synchronized LiDAR capture in a controlled runtime verification setting; it is illustrative only and not a new perception result.}
\label{fig:lidar_runtime_bev}
\end{figure}

\subsection{Sensor rig integration from \texttt{calib\_data.json}}
Perception datasets attach a calibrated sensor suite (multi-camera + LiDAR) to the ego vehicle. The experiments use:
\begin{itemize}
\item camera intrinsics from \texttt{K\_undistortion} (pinhole model),
\item image size from \texttt{image\_size},
\item camera extrinsics from \texttt{cTv} (vehicle$\rightarrow$camera),
\item LiDAR extrinsics from \texttt{vTl} (LiDAR$\rightarrow$vehicle).
\end{itemize}
\paragraph{Camera model and image handling.}
To keep the sensor model consistent across domains, all cameras are evaluated under a pinhole model using $\mathbf{K}_{\text{undist}}$ from \texttt{K\_undistortion}. 
The distortion parameters (\texttt{K}, \texttt{D}) are ignored and no explicit undistortion step is applied to the rendered images; instead, the intrinsic model used for downstream processing is fixed to $\mathbf{K}_{\text{undist}}$ for all runs. 
This prevents calibration modeling choices from confounding domain-gap measurements.



\paragraph{Ego-mounted sensor suite.}
In CARLA, each camera or LiDAR is implemented as an actor that is attached to the ego vehicle using a rigid \texttt{Transform}. \cite{carla_docs_sensors}
The pipeline reads \texttt{calib\_data.json} and instantiates the full sensor suite consistently for every map, which makes downstream comparisons meaningful.

\paragraph{Intrinsics and undistorted pinhole model.}
Although the calibration file provides both camera intrinsics $K$ and distortion coefficients $D$, this thesis uses the undistorted intrinsics $K_{\text{undist}}$ and an ideal pinhole projection model (no distortion).
This avoids simulator-specific distortion implementations and matches the common practice of using rectified imagery as perception input.

\paragraph{Extrinsics and coordinate conventions.}
Extrinsic calibration is represented via 4$\times$4 homogeneous transforms (vehicle $\leftrightarrow$ sensor).
CARLA's coordinate conventions (left-handed Unreal Engine frame and actor-relative transforms) are followed when converting calibration matrices into CARLA sensor poses. \cite{carla_docs_transforms}
Consequently, all camera and LiDAR frames remain consistent across the automatically generated maps and the manually modeled reference map.


If any extrinsics default to the identity transform due to missing data, the run is labeled ``uncalibrated'' because incorrect sensor placement can bias perception results.

\paragraph{Semantic segmentation sensor.}
In addition to RGB cameras and LiDAR, the pipeline optionally attaches a
\texttt{sensor.camera.semantic\_segmentation} actor co-located with each RGB camera,
sharing identical intrinsics and extrinsic attachment transform.
The semantic camera produces pixel-level ground-truth labels using CARLA's built-in
Cityscapes-compatible class encoding, enabling generation of weakly labelled segmentation
masks without manual annotation cost.
Captured semantic frames are saved under a dedicated \texttt{semseg\_raw/} subdirectory
and are tracked independently in the run manifest (\texttt{semseg\_files} counter).
These labels are not used in the primary structural domain-gap analysis but are available
for downstream perception training, domain-adaptation experiments, and sim-to-real transfer
evaluation.

\subsection{Dataset export and run manifests}
For each simulator-valid tile, the pipeline exports synchronized sensor streams and a run manifest containing:
\begin{itemize}
\item tile identifiers and bounds,
\item sensor configuration hash,
\item weather and traffic configuration,
\item random seeds,
\item file lists and checksums.
\end{itemize}
This supports HPC-scale training and evaluation without requiring CARLA-in-the-loop execution on the cluster.

\subsection{Summary}
Simulator-based QA complements offline validation by testing CARLA-native loadability and basic drivability. Dataset generation is standardized via sensor configuration files and manifests to enable reproducible training and cross-domain evaluation.


\paragraph{Crash-resilient tile validation.}
During repeated OpenDRIVE imports on Windows, the CARLA process may terminate abruptly at the engine level (e.g., with exit code \texttt{0xC0000409}), which bypasses Python exception handling.
To prevent a single failing tile from aborting complete experiments, simulator-dependent validation is executed in isolated worker subprocesses (one worker per tile).
If the worker crashes, the main pipeline continues, records the tile outcome, and proceeds with the next tile.
This design turns simulator instability into a measurable outcome (tile pass rate and failure taxonomy) rather than a pipeline-ending event.

\subsection{Determinism audit}
\label{sec:determinism_audit}
\textbf{Question:} Does the automated map generation pipeline produce consistent structural outputs under identical inputs?

\textbf{Procedure:} The OSM-to-OpenDRIVE conversion process is executed multiple times using identical pinned OSM inputs and geographic bounds.
Generated OpenDRIVE files are compared using both cryptographic hashes and semantic structural signatures, including road count, junction count, and lane distribution.
File-level hash differences are analyzed alongside semantic equivalence to distinguish superficial formatting variability from meaningful structural nondeterminism.

\textbf{Expected result:} While file-level hashes may differ due to serialization or ordering effects, semantic structural signatures are expected to remain invariant across runs.
Any deviations are explicitly reported as nondeterministic behavior.

\subsection{Structural domain gap between manual and generated maps}
\label{sec:structural_gap_experiment}
Structural differences between the manually authored map baseline and an automatically generated map are quantified using whole-map statistics and tile-level metrics.
Whole-map metrics include road and lane counts, intersection density, and geometric properties such as curvature distributions.
Tile-level structural differences are aggregated spatially to identify regions with elevated discrepancies.

The primary outputs are (i) a table of global structural statistics and (ii) a tile-level discrepancy heatmap that localizes structural gap intensity across the city.

\subsection{Experiment 0: correspondence validation and frame diagnosis}
\label{sec:exp0_correspondence}
Before any tile-level domain gap is interpreted, correspondence is validated using a three-step protocol:
\begin{enumerate}
  \item \textbf{Identity test (A vs.\ A):} verifies harness correctness. Expected outcome: match rate $\approx 100\%$, distance percentiles near zero, IoU near one.
  \item \textbf{Auto vs.\ auto:} establishes a baseline for origin drift and natural variability under identical coordinate conventions.
  \item \textbf{Manual vs.\ auto:} diagnoses cross-domain compatibility. If a robust translation estimate collapses distance percentiles and IoU gating succeeds, tile-level comparison proceeds on ``good matches'' only; otherwise tile-level analysis is refused.
\end{enumerate}
This protocol ensures that tile-level metrics quantify structural discrepancies rather than mismatched spatial locations.

\section{Perception Capability Results}
\label{sec:perception_results}
\paragraph{Summary of quantitative results.}
RQ1 was assessed over five determinism-audit runs: the verdict is \textbf{NONDETERMINISTIC} (binary hashes diverge; topological counts invariant, $CV=0.0$). RQ2 structural gap is reported from the governed run\_11 raw basis: whole-network RMSE $= 54.84\,$m, matched RMSE $= 24.47\,$m for 73.0\% of sampled points within a 50 m criterion, and Hausdorff distance $= 1663\,$m. These whole-network metrics must be read together with the 4.88$\times$ coverage mismatch (258.6 km auto road length vs.\ 52.9 km manual road length); historical pre-alignment values such as 736.65 m are retained only as workflow context.


\subsection{Reference Rig Validation (Town10HD)}
To validate the sensor rig configuration independent of map generation, a controlled stage probe was performed on the built-in \texttt{Town10HD\_Opt} map.
The probe confirmed correct world identity, ego spawn, and RGB callback success on the built-in map, providing a proxy baseline for rig attachment and callback readiness rather than paired Ingolstadt perceptual evidence.
The stable governed evidence is the Town10 stage-probe artifact in the Scenario B evidence bundle.

\subsection{Generated Map Perception Readiness}
The authoritative governed generated-map attempt is the strict-evidence bundle configured for 50 frames.
However, the final governed status did not close successfully: the strict-evidence \texttt{perception\_status.json} reports \texttt{ok=false} and \texttt{frames\_recorded=0}, while the associated \texttt{pair\_manifest.json} records the 50-frame target.
This supports only a bounded statement that the generated-map arm reached governed strict-evidence execution, not that a submission-grade perception evidence pack was completed.

Notably, partial sensor output was produced within the strict evidence bundle:
\texttt{perception\_status.json} records 17 semantic segmentation frames
(\texttt{semseg\_files: 17}) and 10 LiDAR point cloud files (\texttt{ply\_files: 10}),
alongside 35 PNG files, indicating that sensor attachment and partial data capture
succeeded.
The reported failure mode is \texttt{EMPTY\_RECORDER\_MANIFEST} rather than a total
world-loading or spawn failure, which provides a stronger lower bound on pipeline
capability than a complete failure would imply.

\subsection{Manual Reference Blocker}
Perception capture on the manual reference map (\texttt{Grid0828}) could not be completed. The map loads correctly as a cooked asset (identity confirmed: \texttt{Grid0828/Maps/Grid0828/Grid0828}, 11\,634 spawn points), and all identified repo-local control-flow defects were resolved. In the final governed attempt, the rig-attach path was entered with a valid client but \texttt{spawn\_response\_count}\,=\,0 with persistent streaming transport refusal yielded zero frames, classified as a CARLA runtime boundary during live sensor spawn on the custom cooked map. This runtime boundary precludes a direct image-to-image comparison (RQ3) in this thesis.
